\question[3]
In der Methode \texttt{insert} sind zwei Zeilen vertauscht.
\begin{lstlisting}
    def insert(self, x):
        hilf = Eintrag()
        hilf.inhalt = x
        self.pos.next = hilf
        hilf.next = self.pos.next
\end{lstlisting}
\begin{parts}
\part Welche Zeilen müssen getauscht werden?
\part Was passiert mit dieser falschen Methode, wenn die folgenden Anweisungen ausgeführt werden? Begründe.
\begin{lstlisting}
a = Liste()
a.insert(1)
a.insert(2)
while not a.endpos():
    print(a.elem())
    a.advance()
\end{lstlisting}
\end{parts}
\begin{solutionbox}{3.5cm}
a) Die letzten beiden Zeilen. Der neue Eintrag muss zuerst den Nachfolger übernehmen. \\
b) Die Zahl 2 wird endlos ausgegeben. \texttt{self.pos.next} ist bereits \texttt{hilf}, also zeigt \texttt{hilf.next}
auf \texttt{hilf} selbst. Der neue Eintrag ist sein eigener Nachfolger, \texttt{endpos} wird nie True und der bisherige
Eintrag mit der 1 ist nicht mehr erreichbar.
\end{solutionbox}
